\begin{abstract}
DeepSeek serves its 754B-parameter DeepSeek-V4.1-Flash model with expert parallelism. The experts of each layer are spread over many GPUs, every token is sent to the GPUs that hold its experts, and communication is hidden behind computation. We ask what remains of this design on a single laptop with 24\,GB of memory. Our system, \method{}, keeps the dense weights in memory and leaves the experts on the internal SSD. After routing, it reads each missing expert straight into a preallocated GPU buffer. Decoding a token then becomes a chain in which the router must finish before the reads can start, and the next layer must wait for the experts. The SSD therefore sits idle while the GPU works, and the cluster's overlap techniques recover little of that time. Prefetching the next layer's experts even makes decode slower, because the guesses delay the reads the token is waiting for. We fix this with idle-window speculation, which issues the guesses only after the current layer's reads have landed. This simple change makes decode 11--27\% faster across all cache sizes we tested and leaves the output unchanged. On a MacBook Air, \method{} decodes 1.23 tokens per second while keeping 9\,GB free for macOS, and it stays close to DeepSeek's own deployment on open-ended text. We also show that the model's sparse-attention indexer can use BF16 matrix units without changing its scores, which roughly halves its cost in long prefill.
\end{abstract}
